\chapter{Continuous Probability Done Properly}

\begin{center}\itshape
Pick a real number uniformly at random in $[0,1]$. It equals $\frac12$ with probability $0$, and so does every\\
other number; yet it lands somewhere. Probabilities can no longer be added point by point. What replaces\\
the sum, and is there still a probability for \emph{every} set of numbers?
\end{center}

Chapters~1 to~4 lived on countable sets, where an expectation is a sum. A uniform point of $[0,1]$, a Gaussian
measurement or the time until a radioactive atom decays do not: each single value has probability $0$, and an
expectation must become an integral. This chapter builds that integral, Lebesgue's, from the axioms of Chapter~1
and proves the three theorems that make it powerful: monotone convergence, Fatou's lemma and dominated convergence.
With them, densities, changes of variables, independence and the Gaussian law fit into one framework, and the
proofs of the rest of the book become honest.

The first section explains why a probability cannot be defined on all subsets of $[0,1]$, which is why the axioms
name a $\sigma$-algebra. The second gives the one tool that decides when two probabilities are equal: Dynkin's
$\pi$-$\lambda$ theorem. Then come random variables and their laws, the integral, Lebesgue measure and densities,
changes of variables, independence, and finally polar coordinates, which give the Gaussian integral again and the
Box--Muller method that computers use to draw Gaussian numbers.

\section{Why not every set?}

A uniform random point of $[0,1[$ should give each interval its length, and should not care if everything is
shifted by $t$, wrapping around: write $x\oplus t$ for the fractional part of $x+t$ and $A\oplus t=\{a\oplus t:a\in A\}$.

\begin{theorem}[Vitali, 1905]\label{thm:vitali}
There is no map $\mu:\mathcal P([0,1[)\to[0,1]$ that is countably additive, satisfies $\mu([0,1[)=1$, and is invariant under
shifts: $\mu(A\oplus t)=\mu(A)$ for all $A\subseteq[0,1[$ and $t\in[0,1[$.
\end{theorem}
\begin{proof}
Call $x,y\in[0,1[$ equivalent if $x-y\in\Q$. This is an equivalence relation; by the axiom of choice, let $V$ contain exactly
one point of each class. For each rational $r\in[0,1[$, the sets $V\oplus r$ are pairwise disjoint: if $v\oplus r=v'\oplus r'$ with
$v,v'\in V$, then $v-v'\in\Q$, so $v=v'$ and then $r=r'$. And they cover $[0,1[$: any $x$ is equivalent to some $v\in V$, and
$x=v\oplus r$ with $r=(x-v)\oplus0$, the fractional part of $x-v$, a rational in $[0,1[$. So, by countable additivity and invariance,
$1=\sum_r\mu(V\oplus r)=\sum_r\mu(V)$, a sum of countably many equal terms, which is $0$ or $+\infty$. Contradiction.
\end{proof}

So something must give, and it is ``all subsets''. Lebesgue's solution is to define the uniform probability only
on a $\sigma$-algebra of reasonable sets, the Borel sets below, which contains every set that one can describe
without the axiom of choice. In dimension three the situation is even stranger: Banach and Tarski (1924) cut a ball
into five pieces that can be moved rigidly to form two balls of the same size; the pieces are, necessarily,
not Borel sets.

\section{\texorpdfstring{$\sigma$}{Sigma}-algebras, measures and Dynkin's theorem}

\begin{definition}
A \emph{measurable space} is a set $E$ with a $\sigma$-algebra $\mathcal E$ (Chapter~1). A \emph{measure} on it is a map
$\mu:\mathcal E\to[0,+\infty]$ with $\mu(\varnothing)=0$ that is countably additive. It is \emph{finite} if $\mu(E)<\infty$, a probability if $\mu(E)=1$, and
\emph{$\sigma$-finite} if $E=\bigcup_nE_n$ with $E_n\in\mathcal E$ and $\mu(E_n)<\infty$.
\end{definition}

The rules of Proposition~\ref{prop:rules} hold for any measure, with the same proofs, except that $P(A^c)=1-P(A)$
needs $\mu(E)<\infty$, and continuity along a decreasing sequence $A_1\supseteq A_2\supseteq\cdots$ needs $\mu(A_1)<\infty$.

An intersection of $\sigma$-algebras on $E$ is a $\sigma$-algebra (each axiom passes to intersections). So for any family
$\mathcal C$ of subsets there is a smallest $\sigma$-algebra containing it, the intersection of all of them: the $\sigma$-algebra
\emph{generated} by $\mathcal C$, written $\sigma(\mathcal C)$. The \emph{Borel} $\sigma$-algebra $\mathcal B(\R)$ is generated by the half-lines
$]-\infty,t]$, $t\in\R$; it contains all intervals, since for instance $]a,b]=\,]-\infty,b]\setminus\,]-\infty,a]$ and
$]a,b[\,=\bigcup_n\,]a,b-\frac1n]$, hence all open sets, each of which is a countable union of open intervals with
rational endpoints. Likewise $\mathcal B(\R^d)$ is generated by the boxes $\prod_i\,]-\infty,t_i]$, and contains all open sets of
$\R^d$ (countable unions of open boxes with rational corners).

To prove that two probabilities coincide on a $\sigma$-algebra, it is enough to check it on a well-chosen family.

\begin{definition}
A family $\mathcal C$ of subsets of $E$ is a \emph{$\pi$-system} if $A\cap B\in\mathcal C$ for $A,B\in\mathcal C$. It is a \emph{$\lambda$-system} if $E\in\mathcal C$,
$B\setminus A\in\mathcal C$ whenever $A\subseteq B$ are in $\mathcal C$, and $\bigcup_nA_n\in\mathcal C$ whenever $A_1\subseteq A_2\subseteq\cdots$ are in $\mathcal C$.
\end{definition}

\begin{theorem}[Dynkin's $\pi$-$\lambda$ theorem]\label{thm:dynkin}
If a $\lambda$-system $\mathcal L$ contains a $\pi$-system $\mathcal C$, then it contains $\sigma(\mathcal C)$.
\end{theorem}
\begin{proof}
An intersection of $\lambda$-systems is a $\lambda$-system, so there is a smallest $\lambda$-system $\mathcal D$ containing $\mathcal C$, and $\mathcal D\subseteq\mathcal L$.
We show that $\mathcal D$ is a $\pi$-system; then it is a $\sigma$-algebra (complements are differences $E\setminus A$, finite unions
are complements of intersections of complements, and a countable union is the increasing union of finite ones),
so $\sigma(\mathcal C)\subseteq\mathcal D\subseteq\mathcal L$.

Let $\mathcal D_1=\{A\in\mathcal D:A\cap C\in\mathcal D\text{ for all }C\in\mathcal C\}$. It contains $\mathcal C$, and it is a $\lambda$-system: $E\cap C=C$; if $A\subseteq B$ are in
$\mathcal D_1$ then $(B\setminus A)\cap C=(B\cap C)\setminus(A\cap C)$ with $A\cap C\subseteq B\cap C$; and $(\bigcup A_n)\cap C=\bigcup(A_n\cap C)$, an increasing union. So
$\mathcal D_1=\mathcal D$. Now let $\mathcal D_2=\{A\in\mathcal D:A\cap B\in\mathcal D\text{ for all }B\in\mathcal D\}$. Since $\mathcal D_1=\mathcal D$, $\mathcal D_2$ contains $\mathcal C$, and it is a
$\lambda$-system by the same three computations. So $\mathcal D_2=\mathcal D$, which says that $\mathcal D$ is a $\pi$-system.
\end{proof}

\begin{theorem}[uniqueness of measures]\label{thm:uniqueness}
Let $\mu$, $\nu$ be measures on $\sigma(\mathcal C)$, where $\mathcal C$ is a $\pi$-system, that agree on $\mathcal C$. Suppose that $E=\bigcup_nE_n$ for an
increasing sequence $E_n\in\mathcal C$ with $\mu(E_n)<\infty$. Then $\mu=\nu$. In particular, two probabilities that agree on $\mathcal C$ agree
on $\sigma(\mathcal C)$, if $E\in\mathcal C$.
\end{theorem}
\begin{proof}
Fix $n$. The family $\mathcal L_n$ of sets $A\in\sigma(\mathcal C)$ with $\mu(A\cap E_n)=\nu(A\cap E_n)$ is a $\lambda$-system: it contains $E$; for $A\subseteq B$ in $\mathcal L_n$,
$\mu((B\setminus A)\cap E_n)=\mu(B\cap E_n)-\mu(A\cap E_n)$, a difference of finite numbers, and the same for $\nu$; and increasing unions pass
to the limit by continuity. It contains $\mathcal C$, since $C\cap E_n\in\mathcal C$. By Dynkin, $\mathcal L_n=\sigma(\mathcal C)$. Let $n\to\infty$: by continuity,
$\mu(A)=\lim\mu(A\cap E_n)=\lim\nu(A\cap E_n)=\nu(A)$.
\end{proof}

Note that the half-lines $]-\infty,t]$ form a $\pi$-system, as do the boxes $\prod_i\,]-\infty,t_i]$. Here are the facts that
this tome admits.

\begin{reminder}[Admitted (proved in Tome~6)]
\begin{enumerate}[label=(A\arabic*)]
\item \emph{Existence.} There is a measure $\lambda_d$ on $\mathcal B(\R^d)$, \emph{Lebesgue measure}, giving each box $\prod_i\,]a_i,b_i]$ its volume $\prod_i(b_i-a_i)$;
we write $\lambda=\lambda_1$. For probabilities $\mu_1,\mu_2,\ldots$ on $\mathcal B(\R)$ there is a probability space carrying independent
random variables $X_1,X_2,\ldots$ (Section~\ref{sec:independence}) with laws $\mu_1,\mu_2,\ldots$.
\item \emph{Tonelli and Fubini.} For $\sigma$-finite measures $\mu$ on $(E,\mathcal E)$ and $\nu$ on $(F,\mathcal F)$ there is a measure $\mu\otimes\nu$ on the
$\sigma$-algebra generated by the sets $A\times B$, with $(\mu\otimes\nu)(A\times B)=\mu(A)\nu(B)$, and $\lambda_2=\lambda\otimes\lambda$ on $\mathcal B(\R^2)$. If
$h\ge0$ is measurable on $E\times F$, then $x\mapsto\int h(x,y)\dd\nu(y)$ is measurable and
\[
\begin{aligned}\int h\dd(\mu\otimes\nu)&=\int\Bigl(\int h\dd\nu\Bigr)\dd\mu\eqbr=\int\Bigl(\int h\dd\mu\Bigr)\dd\nu,\end{aligned}
\]
where the inner integrals are in $y$ and in $x$ respectively.
The same holds for $h$ integrable with respect to $\mu\otimes\nu$ (Fubini).
\end{enumerate}
\end{reminder}

The integrals in (A2) are defined in Section~\ref{sec:integral}. By Theorem~\ref{thm:uniqueness}, $\lambda_d$ is unique (the boxes $\prod]-n,n]$
exhaust $\R^d$), and so is the law of the sequence in (A1) on the events that involve finitely many terms.

\section{Random variables and their laws}

\begin{definition}
A map $f:(E,\mathcal E)\to(F,\mathcal F)$ is \emph{measurable} if $f^{-1}(B)\in\mathcal E$ for every $B\in\mathcal F$. A \emph{random variable} is a measurable map
$X:(\Omega,\mathcal F)\to(\R,\mathcal B(\R))$ (or $\R^d$). Its \emph{law} is the probability $\mu_X(B)=P(X\in B)$ on $\mathcal B(\R)$, and its
\emph{distribution function} is $F_X(t)=P(X\le t)$.
\end{definition}

\begin{proposition}\label{prop:measurable}
\begin{enumerate}[label=(\alph*)]
\item If $\mathcal F=\sigma(\mathcal C)$, then $f$ is measurable as soon as $f^{-1}(C)\in\mathcal E$ for all $C\in\mathcal C$. So $X:\Omega\to\R$ is a random variable iff
$\{X\le t\}$ is an event for every $t$, and $X=(X_1,\dots,X_d)$ is a random vector iff each $X_i$ is a random variable.
\item Continuous maps $\R^d\to\R^k$ are measurable (``Borel''). Hence sums, products and continuous functions of random
variables are random variables.
\item If $X_1,X_2,\ldots$ are random variables, so are $\sup_nX_n$, $\inf_nX_n$, $\limsup_nX_n$ and $\liminf_nX_n$ (where finite), and the set
where $(X_n)$ converges is an event.
\item The law of $X$ is determined by $F_X$, which is non-decreasing, right-continuous, with limits $0$ at $-\infty$ and $1$ at $+\infty$.
\end{enumerate}
\end{proposition}
\begin{proof}
(a) The sets $B\subseteq F$ with $f^{-1}(B)\in\mathcal E$ form a $\sigma$-algebra, since preimages commute with complements and countable
unions; it contains $\mathcal C$, hence $\sigma(\mathcal C)$. For a vector, $\{X\in\prod_i]-\infty,t_i]\}=\bigcap_i\{X_i\le t_i\}$.
(b) The preimage of an open set by a continuous map is open, so by (a) it suffices that the open sets generate
$\mathcal B(\R^k)$: they belong to it, and each generating box $\prod_i\,]-\infty,t_i]$ is the complement of an open set. If $X$ and $Y$ are random variables, $(X,Y)$ is a random
vector by (a), and $X+Y$, $XY$ are continuous functions of it; a composition of measurable maps is measurable.
(c) $\{\sup_nX_n\le t\}=\bigcap_n\{X_n\le t\}$, and $\inf_nX_n=-\sup_n(-X_n)$; then $\limsup_n X_n=\inf_n\sup_{k\ge n}X_k$. The set of convergence is
$\{\limsup X_n=\liminf X_n\}$, where both are finite.
(d) The half-lines form a $\pi$-system generating $\mathcal B(\R)$, and $]-\infty,n]$ increases to $\R$, so Theorem~\ref{thm:uniqueness} applies to two
laws with the same distribution function. Monotonicity is clear; right-continuity and the limits follow from continuity of $P$ along
$]-\infty,t+\frac1n]\downarrow\,]-\infty,t]$, $]-\infty,-n]\downarrow\varnothing$ and $]-\infty,n]\uparrow\R$.
\end{proof}

A random variable of Chapter~2 is one of these: if $X$ takes countably many values, $\{X\le t\}$ is a countable union of the
events $\{X=x\}$, $x\le t$.

\section{The integral}\label{sec:integral}

Fix a measure space $(E,\mathcal E,\mu)$; the case to keep in mind is a probability space, where the integral will be the
expectation. Functions take values in $[0,+\infty]$ or in $\R$, and are measurable.

\paragraph{Simple functions.} A \emph{simple} function is $s=\sum_{i=1}^ma_i\mathbf 1_{A_i}$ with $a_i\ge0$ and $A_i\in\mathcal E$. Its integral is
$\int s\dd\mu=\sum_ia_i\mu(A_i)$ (with $0\cdot\infty=0$). This does not depend on the way $s$ is written: the sets $A_i$ cut $E$ into at
most $2^m$ disjoint \emph{atoms} $C_j$ (intersections of $A_i$ or $A_i^c$ for each $i$), each contained in or disjoint from each $A_i$, and
by additivity
\[
\begin{aligned}\sum_ia_i\mu(A_i)&=\sum_ia_i\sum_{j:\,C_j\subseteq A_i}\mu(C_j)\eqbr=\sum_j\Bigl(\sum_{i:\,C_j\subseteq A_i}a_i\Bigr)\mu(C_j)\eqbr=\sum_js(C_j)\,\mu(C_j),\end{aligned}
\]
where $s(C_j)$ is the constant value of $s$ on $C_j$; the right-hand side only depends on $s$ (merge the atoms on which $s$ takes the
same value). Writing two simple functions over common atoms shows that $\int(as+bt)=a\int s+b\int t$ for $a,b\ge0$, and that
$s\le t$ implies $\int s\le\int t$.

\begin{definition}
For measurable $f:E\to[0,+\infty]$, $\displaystyle\int f\dd\mu=\sup\Bigl\{\int s\dd\mu:s\text{ simple},\ 0\le s\le f\Bigr\}\in[0,+\infty]$.
\end{definition}

For simple $f$ this agrees with the previous definition, by monotonicity, and $0\le f\le g$ implies $\int f\le\int g$.

\begin{theorem}[monotone convergence]\label{thm:mct}
If $0\le f_1\le f_2\le\cdots$ are measurable and $f_n\to f$ pointwise, then $\int f_n\dd\mu\to\int f\dd\mu$.
\end{theorem}
\begin{proof}
The sequence $\int f_n$ is non-decreasing and bounded by $\int f$, so its limit $\ell$ satisfies $\ell\le\int f$. Conversely, let
$s=\sum_ia_i\mathbf 1_{B_i}\le f$ be simple, written over disjoint $B_i$, and $0<c<1$. The sets $A_n=\{f_n\ge cs\}$ are measurable and increase
to $E$: if $s(x)=0$ then $x\in A_1$, and otherwise $f_n(x)\to f(x)\ge s(x)>cs(x)$. Then $f_n\ge cs\mathbf 1_{A_n}$, so
\[
\begin{aligned}\int f_n\dd\mu&\ge c\sum_ia_i\mu(B_i\cap A_n)\eqbr\xrightarrow[n\to\infty]{}c\sum_ia_i\mu(B_i)=c\int s\dd\mu\end{aligned}
\]
by continuity of $\mu$ along $B_i\cap A_n\uparrow B_i$. So $\ell\ge c\int s$ for all $c<1$ and all such $s$: $\ell\ge\int f$.
\end{proof}

Every measurable $f\ge0$ is the increasing limit of the simple functions
\[
\varphi_n(f)=\min\Bigl(n,\;2^{-n}\lfloor2^nf\rfloor\Bigr),
\]
which take finitely many values on measurable sets ($\varphi_n(f)$ increases with $n$ because $2^{-n}\lfloor2^nx\rfloor\le2^{-n-1}\lfloor2^{n+1}x\rfloor$). With
monotone convergence this transfers the properties of simple functions to all non-negative ones.

\begin{figure}[ht]
\centering
\fig{ch05-simple}
\caption{The simple functions $\varphi_1(f)$ (orange) and $\varphi_3(f)$ (blue) below a function $f$ (black). Lebesgue cuts the \emph{values} into
slices of height $2^{-n}$, where Riemann cuts the \emph{axis} into intervals; the sets where $f$ lies in a slice may be complicated, but
they only need to have a measure.}
\label{fig:simple}
\end{figure}

\begin{corollary}\label{cor:linearity}
For measurable $f,g,f_n\ge0$ and $a,b\ge0$: $\int(af+bg)=a\int f+b\int g$, and $\int\sum_nf_n=\sum_n\int f_n$.
\end{corollary}
\begin{proof}
$a\varphi_n(f)+b\varphi_n(g)$ is simple and increases to $af+bg$; apply monotone convergence to both sides of the identity for simple
functions. For the series, apply this to the partial sums, which increase to $\sum_nf_n$, and use monotone convergence again.
\end{proof}

\begin{theorem}[Fatou's lemma]\label{thm:fatou}
For measurable $f_n\ge0$, $\displaystyle\int\liminf_nf_n\dd\mu\le\liminf_n\int f_n\dd\mu$.
\end{theorem}
\begin{proof}
$g_n=\inf_{k\ge n}f_k$ increases to $\liminf f_n$, and $g_n\le f_k$ for $k\ge n$, so $\int g_n\le\inf_{k\ge n}\int f_k$. Let $n\to\infty$ and apply monotone
convergence on the left.
\end{proof}

\begin{definition}
A measurable $f:E\to\R$ is \emph{integrable} if $\int|f|\dd\mu<\infty$; then $\int f\dd\mu=\int f^+\dd\mu-\int f^-\dd\mu$, where $f^+=\max(f,0)$ and
$f^-=\max(-f,0)$. On a probability space, $\mathbb E[X]=\int X\dd P$.
\end{definition}

The integral is linear and monotone on integrable functions, and $|\int f|\le\int|f|$. For linearity, if $h=f+g$ then
$h^+-h^-=f^+-f^-+g^+-g^-$, so $h^++f^-+g^-=h^-+f^++g^+$; integrate both sides with Corollary~\ref{cor:linearity} and rearrange the finite
numbers. Multiplication by a constant is direct, and $\int f\le\int g$ for $f\le g$ follows from $\int(g-f)\ge0$.

\begin{theorem}[dominated convergence]\label{thm:dct}
Let $f_n\to f$ pointwise, with $|f_n|\le g$ for all $n$ and $g$ integrable. Then $f$ is integrable and $\int f_n\dd\mu\to\int f\dd\mu$; in fact
$\int|f_n-f|\dd\mu\to0$.
\end{theorem}
\begin{proof}
$|f|\le g$, so $f$ is integrable. The functions $2g-|f_n-f|$ are non-negative and tend to $2g$. By Fatou,
$\int2g\le\liminf_n\int(2g-|f_n-f|)=\int2g-\limsup_n\int|f_n-f|$. Since $\int2g<\infty$, $\limsup\int|f_n-f|\le0$. Finally
$|\int f_n-\int f|\le\int|f_n-f|$.
\end{proof}

\begin{think}
Let $f_n=n\,\mathbf 1_{]0,1/n[}$ on $\R$ with Lebesgue measure. Where do the functions converge, and what are their integrals? Which
hypothesis of each of the three theorems fails?
\end{think}

A property holds \emph{almost everywhere} (a.e.) if it fails only on a set of measure $0$. Changing $f$ on a null set does
not change its integral (a non-negative function that vanishes a.e.\ has integral $0$, since every simple function
below it does; apply this to $f-\min(f,g)$ and $g-\min(f,g)$ when $f=g$ a.e.), so the convergence theorems
hold with ``a.e.'' hypotheses. Markov's inequality $\mu(f\ge a)\le\frac1a\int f$ holds with the proof of Chapter~4; in particular
$\int f=0$ for $f\ge0$ implies $\mu(f\ge\frac1k)=0$ for all $k$, so $f=0$ a.e.

\begin{proposition}\label{prop:consistency}
\begin{enumerate}[label=(\alph*)]
\item If $X\ge0$ takes countably many values, $\mathbb E[X]=\sum_xx\,P(X=x)$, as in Chapter~2.
\item \emph{Transfer}: for a random variable $X$ and a Borel $g\ge0$, $\mathbb E[g(X)]=\int g\dd\mu_X$; the same holds for $g(X)$ integrable.
\end{enumerate}
\end{proposition}
\begin{proof}
(a) $X=\sum_xx\mathbf 1_{\{X=x\}}$, a countable sum of non-negative functions; apply Corollary~\ref{cor:linearity}. (b) For $g=\mathbf 1_B$, both sides are
$P(X\in B)$, by the definition of the law. By linearity it holds for simple $g$, by monotone convergence with $\varphi_n(g)$ for
$g\ge0$, and by splitting $g=g^+-g^-$ in the integrable case.
\end{proof}

All the results of Chapters~2 and~4 that only used linearity and positivity, such as Cauchy--Schwarz, the
variance of a sum, Markov, Chebyshev, Jensen and Hoeffding's lemma, now hold for all random variables, with the same
proofs. The product formula $\mathbb E[XY]=\mathbb E[X]\mathbb E[Y]$ for independent variables needs a new proof, given in
Section~\ref{sec:independence}; with it, the weak law, Chernoff and Hoeffding hold in general too.

\section{Lebesgue measure and densities}

We write $\int_Bf(x)\dd x$ for $\int f\mathbf 1_B\dd\lambda$. The first task is to check that this is the integral of Tome~0 when both make sense.

\begin{proposition}\label{prop:riemann}
\begin{enumerate}[label=(\alph*)]
\item If $f\ge0$ is continuous on $[a,b]$, then $\int_{[a,b]}f\dd\lambda=\int_a^bf$, the integral of Tome~0.
\item If $f\ge0$ is continuous on an open interval $]a,b[$ ($-\infty\le a<b\le+\infty$), then $\int_{]a,b[}f\dd\lambda=\lim\int_{a_n}^{b_n}f$ for any $a_n\downarrow a$,
$b_n\uparrow b$; in particular the improper integral of Tome~0 is a Lebesgue integral.
\end{enumerate}
\end{proposition}
\begin{proof}
(a) Cut $[a,b]$ into $2^n$ intervals $I_{n,k}$ of length $\frac{b-a}{2^n}$, half-open except the last, and let $m_{n,k}=\min_{I_{n,k}}f$ (minimum over the
closed interval). The step function $s_n=\sum_km_{n,k}\mathbf 1_{I_{n,k}}$ is simple, with $\int s_n\dd\lambda=\sum_km_{n,k}\frac{b-a}{2^n}=:L_n$, and $s_n$ increases with $n$
(the minimum over a sub-interval is at least the minimum over the interval). Given $\varepsilon>0$, Heine's Lemma~\ref{lem:heine} (after
the affine change $[0,1]\to[a,b]$) gives $\delta$ with $|f(x)-f(y)|<\varepsilon$ for $|x-y|<\delta$, so $0\le f-s_n\le\varepsilon$ on $[a,b]$ as soon as
$\frac{b-a}{2^n}<\delta$: $s_n\to f$. By monotone convergence, $\int_{[a,b]}f\dd\lambda=\lim L_n$. On the other side, the monotonicity and additivity of the
integral of Tome~0 give $L_n\le\int_a^bf\le L_n+\varepsilon(b-a)$ for such $n$. So both are equal to $\lim L_n$.

(b) $f\mathbf 1_{[a_n,b_n]}$ increases to $f\mathbf 1_{]a,b[}$; apply (a) and monotone convergence.
\end{proof}

With Tome~0, Exercise~6.7 and the change $x=\frac u{\sqrt2}$, $\int_\R e^{-u^2/2}\dd u=2\sqrt2\int_0^\infty e^{-x^2}\dd x=\sqrt{2\pi}$.

\begin{think}
A dart lands uniformly at random in $[0,1]$. What is the probability that it hits a rational number? That the decimal
expansion of the point contains no digit $7$? Both sets are infinite, and the second is uncountable.
\end{think}

\begin{definition}
A Borel function $f\ge0$ on $\R$ with $\int_\R f\dd\lambda=1$ is a \emph{probability density}. A random variable $X$ \emph{has density} $f$ if
$P(X\in B)=\int_Bf(x)\dd x$ for every Borel set $B$.
\end{definition}

By the uniqueness theorem it is enough that $P(X\le t)=\int_{-\infty}^tf(x)\dd x$ for all $t$. So if $F_X$ is continuous, and of class $C^1$ outside
finitely many points, then $F_X'$ (defined anywhere at those points) is a density of $X$, by the fundamental theorem of
calculus on each interval between the exceptional points. A density is unique up to a null set (Exercise~\ref{ex:densityunique}),
and by transfer (Proposition~\ref{prop:consistency}(b), then simple functions and monotone convergence),
\[
\mathbb E[g(X)]=\int_\R g(x)f(x)\dd x
\]
for Borel $g\ge0$, and for Borel $g$ with $g(X)$ integrable.

Back to the dart: a single point has $\lambda$-measure $0$, since it lies in intervals of arbitrarily small length, so by countable
additivity every countable set, such as $\Q\cap[0,1]$, has measure $0$, and the dart never hits a rational. The points of $[0,1[$
with no $7$ among their first $n$ decimals form $9^n$ intervals of length $10^{-n}$, so the probability of never seeing a $7$ is at most
$0.9^n$ for every $n$: it is $0$, although the set is uncountable.

\begin{example}\label{ex:laws}
\begin{itemize}
\item \emph{Uniform} on $[a,b]$: density $\frac1{b-a}\mathbf 1_{[a,b]}$. Mean $\frac{a+b}2$.
\item \emph{Exponential}$(\theta)$, $\theta>0$: density $\theta e^{-\theta x}$ on $]0,\infty[$, $P(X>t)=e^{-\theta t}$, mean $\frac1\theta$, variance $\frac1{\theta^2}$.
\item \emph{Normal} $N(m,\sigma^2)$, with $\sigma>0$: the density is $\frac1{\sigma\sqrt{2\pi}}e^{-(x-m)^2/(2\sigma^2)}$, the mean $m$, the variance $\sigma^2$. For $N(0,1)$, integrating $x\cdot xe^{-x^2/2}$ by parts
gives $\mathbb E[X^2]=1$.
\item \emph{Gamma}$(a)$, $a>0$: density $\frac1{\Gamma(a)}x^{a-1}e^{-x}$ on $]0,\infty[$, where $\Gamma(a)=\int_0^\infty x^{a-1}e^{-x}\dd x$ (finite: compare with
$x^{a-1}$ near $0$ and with $e^{-x/2}$ near $\infty$). Exponential$(1)$ is Gamma$(1)$.
\item \emph{Cauchy}: density $\frac1{\pi(1+x^2)}$, which integrates to $\frac1\pi[\arctan x]_{-\infty}^{+\infty}=1$.
\end{itemize}
\end{example}

\begin{proposition}[inverse transform]\label{prop:inverse}
Let $F$ be non-decreasing and right-continuous, with limits $0$ and $1$ at $\mp\infty$, and $G(u)=\inf\{x:F(x)\ge u\}$ for $u\in\,]0,1[$. If $U$ is
uniform on $]0,1[$, then $G(U)$ has distribution function $F$. So every such $F$ is a distribution function.
\end{proposition}
\begin{proof}
For $u\in\,]0,1[$ the set $\{x:F(x)\ge u\}$ is non-empty (limit $1$), bounded below (limit $0$), and closed on the left by right-continuity:
it is $[G(u),+\infty[$. So $G(u)\le x$ iff $F(x)\ge u$. Then $G$ is non-decreasing, hence Borel ($\{G\le x\}$ is an interval), and
$P(G(U)\le x)=P(U\le F(x))=F(x)$.
\end{proof}

For instance $-\frac1\theta\ln U$ is exponential$(\theta)$, and $\tan(\pi(U-\frac12))$ is Cauchy. This is how computers turn uniform random numbers
into any law whose $G$ is computable; for the Gaussian law, $G$ has no closed form and Section~\ref{sec:polar} does better.

\section{Changing variables}

\begin{proposition}\label{prop:cov}
Let $I$, $J$ be open intervals and $\varphi:J\to I$ a bijection of class $C^1$ with $\varphi'\ne0$. For every Borel $h\ge0$ on $I$,
\[
\int_Ih(x)\dd x=\int_Jh(\varphi(u))\,|\varphi'(u)|\dd u.
\]
\end{proposition}
\begin{proof}
For a Borel set $B\subseteq I$, let
\[
\nu(B)=\int_J\mathbf 1_B(\varphi(u))\,|\varphi'(u)|\dd u;
\]
the integrand is Borel, since $\varphi$ is continuous. By
Corollary~\ref{cor:linearity}, $\nu$ is countably additive: a measure on $I$. Since $\varphi'$ is continuous and does not vanish, it has
constant sign and $\varphi$ is strictly monotone. For $[c,d]\subseteq I$, $\varphi^{-1}([c,d])$ is a segment $[\alpha,\beta]\subseteq J$, and by
Proposition~\ref{prop:riemann} and the fundamental theorem,
\[
\begin{aligned}\nu(]c,d])&=\int_\alpha^\beta|\varphi'|=|\varphi(\beta)-\varphi(\alpha)|\eqbr=d-c=\lambda(]c,d])\end{aligned}
\]
(single points have measure $0$ for both: $\nu(\{c\})=\int\mathbf 1_{\{\varphi^{-1}(c)\}}|\varphi'|=0$). The intervals $]c,d]$ with $[c,d]\subseteq I$, together with $\varnothing$, form a
$\pi$-system that generates the Borel sets of $I$, and an increasing sequence of them exhausts $I$. By Theorem~\ref{thm:uniqueness},
$\nu=\lambda$ on $I$: the formula holds for indicators. By linearity it holds for simple $h$, and by monotone convergence
with $\varphi_n(h)$ for all $h\ge0$.
\end{proof}

\begin{corollary}[density of $g(X)$]\label{cor:densityg}
Let $X$ have a density $f$ with $P(X\in I)=1$ for an open interval $I$, and let $g:I\to J$ be a bijection of class $C^1$ with $g'\ne0$
onto an open interval $J$. Then $Y=g(X)$ has the density
\[
f_Y(y)=f\bigl(g^{-1}(y)\bigr)\,\bigl|(g^{-1})'(y)\bigr|
\]
for $y\in J$, and $f_Y=0$ outside $J$.
\end{corollary}
\begin{proof}
$g^{-1}:J\to I$ is $C^1$ with non-zero derivative $\frac1{g'\circ g^{-1}}$. For Borel $B$, by transfer and the proposition with $\varphi=g^{-1}$,
\[
\begin{aligned}P(Y\in B)&=\int_I\mathbf 1_B(g(x))f(x)\dd x\eqbr=\int_J\mathbf 1_B(y)\,f(g^{-1}(y))\,|(g^{-1})'(y)|\dd y.\end{aligned}
\]
\end{proof}

For example, if $X$ is $N(0,1)$ then $m+\sigma X$ is $N(m,\sigma^2)$ for $\sigma>0$. And $|X|$ has density $2\varphi$ on $]0,\infty[$, where
$\varphi(x)=\frac1{\sqrt{2\pi}}e^{-x^2/2}$, so $X^2=g(|X|)$ with $g(x)=x^2$ has density $2\varphi(\sqrt y)\frac1{2\sqrt y}=\frac1{\sqrt{2\pi y}}e^{-y/2}$ on $]0,\infty[$. Comparing
with the Gamma density after the change $y=2z$ gives, as a bonus, $\Gamma(\frac12)=\sqrt\pi$ (Exercise~\ref{ex:gamma}).

\section{Independence}\label{sec:independence}

\begin{definition}
Random variables (or vectors) $X_1,\dots,X_n$ are \emph{independent} if $P(X_1\in B_1,\dots,X_n\in B_n)=\prod_iP(X_i\in B_i)$ for all Borel sets
$B_i$; an infinite family is independent if every finite subfamily is.
\end{definition}

\begin{proposition}\label{prop:indepcdf}
\begin{enumerate}[label=(\alph*)]
\item Real random variables $X_1,\dots,X_n$ are independent as soon as $P(X_1\le t_1,\dots,X_n\le t_n)=\prod_iP(X_i\le t_i)$ for all $t_i\in\R$.
\item If $X_1,\dots,X_n$ are independent, so are $g_1(X_1),\dots,g_n(X_n)$ for Borel $g_i$; and so are functions of disjoint groups.
\item If $X$, $Y$ are independent and integrable, then $XY$ is integrable and $\mathbb E[XY]=\mathbb E[X]\,\mathbb E[Y]$.
\end{enumerate}
\end{proposition}
\begin{proof}
(a) Letting some $t_i\to+\infty$ (continuity), the product formula also holds when some of the half-lines are replaced by $\R$.
Fix half-lines $B_2,\dots,B_n$ (or $\R$). Both sides of the product formula, as functions of $B_1$, are finite measures that agree on
the $\pi$-system of half-lines and on $\R$; by Theorem~\ref{thm:uniqueness} they agree for every Borel $B_1$. Now fix a Borel $B_1$ and
half-lines $B_3,\dots,B_n$, and repeat with $B_2$; after $n$ steps all the $B_i$ are arbitrary Borel sets.
(b) $\{g_i(X_i)\in B_i\}=\{X_i\in g_i^{-1}(B_i)\}$. For two groups, let $V=(X_1,\dots,X_k)$ and $W=(X_{k+1},\dots,X_n)$. If $A$ and $B$ are
products of Borel sets, $P(V\in A,W\in B)=P(V\in A)P(W\in B)$ is the definition of independence. Fix such a $B$: both sides are
finite measures in $A$ that agree on the boxes $\prod_i\,]-\infty,t_i]$, a $\pi$-system generating $\mathcal B(\R^k)$ and exhausting $\R^k$, so they agree
for all Borel $A$. Then fix a Borel $A$ and do the same in $B$. So $V$ and $W$ are independent, and so are functions of them;
more groups are handled in the same way.
(c) For simple $X=\sum a_i\mathbf 1_{\{X=a_i\}}$ and $Y=\sum b_j\mathbf 1_{\{Y=b_j\}}$ (with $a_i,b_j\ge0$),
$\mathbb E[XY]=\sum_{i,j}a_ib_jP(X=a_i,Y=b_j)=\mathbb E[X]\,\mathbb E[Y]$. For $X,Y\ge0$, $\varphi_n(X)$ and $\varphi_n(Y)$ are independent by (b), simple, and their
product increases to $XY$; monotone convergence gives the formula in $[0,\infty]$. Applied to $|X|,|Y|$ it gives integrability, and
in general expand $XY=(X^+-X^-)(Y^+-Y^-)$, where $X^\pm$ and $Y^\pm$ are independent by (b).
\end{proof}

The binary digits of a uniform number are the coins of Chapter~1, and conversely: this is the equivalence announced
there.

\begin{proposition}\label{prop:digits}
Let $U$ be uniform on $[0,1[$ and $D_k=\lfloor2^kU\rfloor-2\lfloor2^{k-1}U\rfloor$ its $k$-th binary digit. Then $D_1,D_2,\ldots$ are independent fair bits.
Conversely, if $B_1,B_2,\ldots$ are independent fair bits, $\sum_kB_k2^{-k}$ is uniform on $[0,1]$.
\end{proposition}
\begin{proof}
$D_k$ is a Borel function of $U$, with values in $\{0,1\}$. For bits $b_1,\dots,b_n$, the event $\{D_1=b_1,\dots,D_n=b_n\}$ is $\{U\in[x,x+2^{-n}[\}$
with $x=\sum_{k\le n}b_k2^{-k}$, of probability $2^{-n}$. So the first $n$ digits have the joint law of $n$ independent fair bits, for every $n$; the
product formula for any finite subfamily follows by summing over the digits not in the subfamily.

Conversely let $V=\sum_kB_k2^{-k}$ and $t=j2^{-n}$ with $0\le j\le2^n$. If $M=\sum_{k\le n}B_k2^{n-k}$, then $M2^{-n}\le V\le M2^{-n}+2^{-n}$. So $\{M<j\}\subseteq\{V\le t\}\subseteq\{M<j\}\cup\{V=t\}$,
and $P(V=t)\le P(M=j-1,\ B_k=1\text{ for }n<k\le m)+P(M=j,\ B_k=0\text{ for }n<k\le m)\le2\cdot2^{-(m-n)}$ for every $m>n$, so $P(V=t)=0$.
Hence $P(V\le t)=P(M<j)=j2^{-n}=t$, since $M$ is uniform on $\{0,\dots,2^n-1\}$. A distribution function is right-continuous and the dyadic numbers
are dense, so $P(V\le t)=t$ for all $t\in[0,1]$: $V$ has the law of $U$.
\end{proof}

So (A1) for coin flips is equivalent to the existence of the uniform law on $[0,1]$. With Corollary~\ref{cor:slln}, almost every number
in $[0,1]$ has asymptotically as many $0$s as $1$s in its binary expansion (Borel, 1909).

\begin{proposition}[joint densities]\label{prop:joint}
Let $X$, $Y$ have densities $f$, $g$. They are independent iff $(X,Y)$ has the density $f(x)g(y)$ on $\R^2$, that is,
$P((X,Y)\in B)=\iint_Bf(x)g(y)\dd x\dd y$ for all Borel $B\subseteq\R^2$. If $(X,Y)$ has a density $h$, then $X$ has density $\int_\R h(x,y)\dd y$.
\end{proposition}
\begin{proof}
If $(X,Y)$ has density $fg$, then by Tonelli $P(X\le s,Y\le t)=\int_{-\infty}^sf\int_{-\infty}^tg=P(X\le s)P(Y\le t)$, and Proposition~\ref{prop:indepcdf}(a)
gives independence. Conversely, if they are independent, the probabilities $B\mapsto P((X,Y)\in B)$ and $B\mapsto\iint_Bfg$ agree on the
boxes $]-\infty,s]\times]-\infty,t]$ (by the same Tonelli computation), which form a $\pi$-system generating $\mathcal B(\R^2)$; so they are equal.
For the marginal, $P(X\le s)=\iint\mathbf 1_{x\le s}h=\int_{-\infty}^s\bigl(\int h(x,y)\dd y\bigr)\dd x$ by Tonelli.
\end{proof}

\begin{proposition}[sums of independent Gaussians]\label{prop:gausssum}
If $X$, $Y$ are independent with densities $f$, $g$, then $X+Y$ has the density $(f*g)(z)=\int_\R f(z-y)g(y)\dd y$. If $X$ is
$N(m,\sigma^2)$ and $Y$ is $N(m',\tau^2)$, then $X+Y$ is $N(m+m',\sigma^2+\tau^2)$.
\end{proposition}
\begin{proof}
By the previous proposition and Tonelli, then the change of variables $x=z-y$ in the inner integral (Proposition~\ref{prop:cov} with a
translation), then Tonelli again,
\[
\begin{aligned}
P(X+Y\le t)&=\int_\R g(y)\Bigl(\int_\R\mathbf 1_{x\le t-y}f(x)\dd x\Bigr)\dd y\eqbr=\int_\R g(y)\Bigl(\int_{-\infty}^tf(z-y)\dd z\Bigr)\dd y\\
&=\int_{-\infty}^t(f*g)(z)\dd z.
\end{aligned}
\]
For the Gaussian case, subtracting the means we may take $m=m'=0$. Let $s=\sigma^2+\tau^2$. Completing the square in $y$,
\[
\frac{(z-y)^2}{\sigma^2}+\frac{y^2}{\tau^2}=\frac s{\sigma^2\tau^2}\Bigl(y-\frac{\tau^2z}s\Bigr)^2+\frac{z^2}s.
\]
The integral over $y$ of $e^{-\frac s{2\sigma^2\tau^2}(y-c)^2}$ is $\sqrt{2\pi\sigma^2\tau^2/s}$ (translation and scaling of $\int e^{-u^2/2}=\sqrt{2\pi}$). So
$(f*g)(z)=\frac1{2\pi\sigma\tau}\sqrt{\frac{2\pi\sigma^2\tau^2}s}\,e^{-z^2/(2s)}=\frac1{\sqrt{2\pi s}}e^{-z^2/(2s)}$.
\end{proof}

\section{Polar coordinates and Box--Muller}\label{sec:polar}

\begin{proposition}[polar coordinates]\label{prop:polar}
For every Borel $h\ge0$ on $\R^2$,
\[
\begin{aligned}\iint_{\R^2}h(x,y)\dd x\dd y&=\eqbr\int_0^\infty\!\!\int_0^{2\pi}h(r\cos\theta,r\sin\theta)r\dd\theta\dd r\end{aligned}
\]
in $[0,+\infty]$.
\end{proposition}
\begin{proof}
By Tonelli, $\lambda_2(\R\times\{0\})=\int_\R\lambda(\{0\})\dd x=0$, so we may integrate over the half-planes $y>0$ and $y<0$ separately. On $y>0$ we change variables
three times, each time in a one-dimensional inner integral, using Tonelli in between. For fixed $y>0$, $x=yt$ ($t\in\R$) gives
\[
\begin{aligned}\iint_{y>0}h&=\int_0^\infty\Bigl(\int_\R h(yt,y)\,y\dd t\Bigr)\dd y\eqbr=\int_\R\Bigl(\int_0^\infty h(yt,y)\,y\dd y\Bigr)\dd t.\end{aligned}
\]
For fixed $t$, let $c_t=(1+t^2)^{-1/2}$; the change $y=c_tr$ ($r>0$) turns the inner integral into $\int_0^\infty h(c_trt,c_tr)\,c_t^2r\dd r$. Exchanging again,
\[
\iint_{y>0}h=\int_0^\infty\Bigl(\int_\R h(c_trt,c_tr)\,c_t^2\dd t\Bigr)r\dd r.
\]
Finally $t=\cot\theta$ is a decreasing $C^1$ bijection from $]0,\pi[$ onto $\R$ with $|\frac{\dd t}{\dd\theta}|=\frac1{\sin^2\theta}=1+t^2$, and $\frac t{\sqrt{1+t^2}}=\cos\theta$,
$c_t=\frac1{\sqrt{1+t^2}}=\sin\theta$ (as $\sin\theta>0$), so $c_t^2\,|\frac{\dd t}{\dd\theta}|=1$. So the inner integral is $\int_0^\pi h(r\cos\theta,r\sin\theta)\dd\theta$. On $y<0$, apply this to
$(x,y)\mapsto h(x,-y)$ and substitute $\theta\mapsto2\pi-\theta$, which gives the integral over $]\pi,2\pi[$. Add the two halves.
\end{proof}

\begin{figure}[ht]
\centering
\fig{ch05-polar}
\caption{Polar coordinates. The cell $r\le\rho\le r+\Delta r$, $\theta\le\phi\le\theta+\Delta\theta$ has area $\frac{(r+\Delta r)^2-r^2}2\Delta\theta\approx r\,\Delta r\,\Delta\theta$: cells far
from the origin are larger, which is the factor $r$ in Proposition~\ref{prop:polar}.}
\label{fig:polar}
\end{figure}

With $h(x,y)=e^{-(x^2+y^2)/2}$, Tonelli and the proposition give $\bigl(\int_\R e^{-x^2/2}\dd x\bigr)^2=2\pi\int_0^\infty re^{-r^2/2}\dd r=2\pi$: Poisson's
proof of the Gaussian integral, which Tome~0 avoided, is now complete.

\begin{theorem}[Box--Muller, 1958]\label{thm:boxmuller}
Let $U_1,U_2$ be independent and uniform on $]0,1[$, and
\[
\begin{gathered}R=\sqrt{-2\ln U_1},\qquad\Theta=2\pi U_2,\qq X=R\cos\Theta,\qquad Y=R\sin\Theta.\end{gathered}
\]
Then $X$ and $Y$ are independent $N(0,1)$ variables.
\end{theorem}
\begin{proof}
By Proposition~\ref{prop:joint}, $(U_1,U_2)$ has density $\mathbf 1_{]0,1[^2}$. For a Borel $B\subseteq\R^2$, by Tonelli,
with $T(u,v)=\bigl(\sqrt{-2\ln u}\cos2\pi v,\sqrt{-2\ln u}\sin2\pi v\bigr)$,
\[
P((X,Y)\in B)=\int_0^1\!\!\int_0^1\mathbf 1_B\bigl(T(u,v)\bigr)\dd u\dd v.
\]
In the inner integral put $u=e^{-r^2/2}$, a decreasing $C^1$ bijection from $]0,\infty[$ onto $]0,1[$ with $|\frac{\dd u}{\dd r}|=re^{-r^2/2}$; in the outer one put
$v=\frac\theta{2\pi}$. Then exchange the integrals and use Proposition~\ref{prop:polar} with $h(x,y)=\mathbf 1_B(x,y)\,\varphi(x)\varphi(y)$, where
$\varphi(x)=\frac1{\sqrt{2\pi}}e^{-x^2/2}$, so that $h(r\cos\theta,r\sin\theta)=\mathbf 1_B(r\cos\theta,r\sin\theta)\frac{e^{-r^2/2}}{2\pi}$:
\[
\begin{aligned}P((X,Y)\in B)&=\int_0^\infty\!\!\int_0^{2\pi}h(r\cos\theta,r\sin\theta)\,r\dd\theta\dd r\eqbr=\iint_B\varphi(x)\varphi(y)\dd x\dd y.\end{aligned}
\]
By Proposition~\ref{prop:joint}, $X$ and $Y$ are independent and $N(0,1)$.
\end{proof}

\begin{figure}[ht]
\centering
\fig{ch05-boxmuller}
\caption{Box--Muller. Left: $400$ uniform points in the square, coloured by $U_2$. Right: their images $(X,Y)$, a cloud of independent
Gaussian pairs. The horizontal coordinate $U_1$ becomes the radius (small $U_1$, far out) and $U_2$ the angle.}
\label{fig:boxmuller}
\end{figure}

\begin{beauty}
The Gaussian law has no elementary inverse distribution function, so the inverse transform does not apply directly.
Box and Muller noticed that a \emph{pair} of Gaussians is easy: its law is invariant under rotation, so the angle is
uniform, and its squared radius is exponential. The difficult one-dimensional problem becomes easy in two
dimensions, and Poisson's trick for the Gaussian integral is the same idea, read as a computation rather than as an
algorithm.
\end{beauty}

\section{Exercises}

\begin{exercise}[$\star$]\label{ex:counter}
Let $f_n=n\mathbf 1_{]0,1/n[}$ on $\R$. Show that $f_n\to0$ everywhere while $\int f_n\dd\lambda=1$. Which hypothesis fails in monotone and in
dominated convergence? Show that Fatou's inequality is strict here.
\end{exercise}

\begin{exercise}[$\star$]\label{ex:densityunique}
Show that if $f$ and $g$ are densities of $X$, then $f=g$ almost everywhere.
\end{exercise}

\begin{exercise}[$\star\star$]\label{ex:tail}
Let $X\ge0$ be a random variable. Show that $\mathbb E[X]=\int_0^\infty P(X>t)\dd t$, and that $\mathbb E[X^2]=\int_0^\infty2tP(X>t)\dd t$. Compare with
Theorem~\ref{thm:linearity}(d).
\end{exercise}

\begin{exercise}[$\star\star$]\label{ex:cauchy}
Let $X$, $Y$ be independent $N(0,1)$ variables. (a) Show that $Y/X$ (defined when $X\ne0$, an event of probability $1$) is Cauchy. (b) Show that a
Cauchy variable is not integrable. (c) Show that $(X/R,Y/R)$, with $R=\sqrt{X^2+Y^2}$, is uniform on the unit circle (its angle is
uniform on $[0,2\pi[$) and independent of $R$, and that $R^2$ is exponential with mean $2$.
\end{exercise}

\begin{exercise}[$\star\star$]\label{ex:gamma}
(a) Show that $\Gamma(a+1)=a\Gamma(a)$ for $a>0$, $\Gamma(n)=(n-1)!$ and $\Gamma(\frac12)=\sqrt\pi$. (b) Show that the sum of $n$ independent
exponential$(1)$ variables is Gamma$(n)$.
\end{exercise}

\begin{exercise}[$\star\star$]\label{ex:order}
Let $U_1,\dots,U_n$ be independent uniform on $[0,1]$ and $U_{(k)}$ the $k$-th smallest. Show that $U_{(k)}$ has density
$n\binom{n-1}{k-1}t^{k-1}(1-t)^{n-k}$ on $[0,1]$, and that $\mathbb E[U_{(k)}]=\frac k{n+1}$.
\end{exercise}

\begin{exercise}[$\star\star$]\label{ex:cantor}
Let $B_1,B_2,\ldots$ be independent fair bits and $X=\sum_k2B_k3^{-k}$. Let $C$ be the Cantor set, the numbers of $[0,1]$ that have a
ternary expansion with digits in $\{0,2\}$. Show that $\lambda(C)=0$, that $P(X\in C)=1$ and that $F_X$ is continuous. Deduce that $X$ has no
density, although no single value has positive probability.
\end{exercise}

\begin{exercise}[$\star\star\star$]\label{ex:betagamma}
\emph{The beta--gamma algebra.} Let $X$, $Y$ be independent, Gamma$(a)$ and Gamma$(b)$. Show that $S=X+Y$ and $V=\frac X{X+Y}$ are independent,
that $S$ is Gamma$(a+b)$ and that $V$ has density $\frac{\Gamma(a+b)}{\Gamma(a)\Gamma(b)}v^{a-1}(1-v)^{b-1}$ on $]0,1[$ (the \emph{Beta}$(a,b)$ law). Deduce
$\int_0^1v^{a-1}(1-v)^{b-1}\dd v=\frac{\Gamma(a)\Gamma(b)}{\Gamma(a+b)}$.
\end{exercise}

\section{Hints}

\begin{hint}{ex:counter}
No integrable function dominates all the $f_n$: $\sup_nf_n(x)$ is about $\frac1x$.
\end{hint}
\begin{hint}{ex:densityunique}
Integrate over $B=\{f>g\}$.
\end{hint}
\begin{hint}{ex:tail}
Write $X=\int_0^\infty\mathbf 1_{t<X}\dd t$ and use Tonelli on $\Omega\times[0,\infty[$.
\end{hint}
\begin{hint}{ex:cauchy}
Use the representation of Theorem~\ref{thm:boxmuller}: $Y/X=\tan\Theta$.
\end{hint}
\begin{hint}{ex:gamma}
(a) Integrate by parts; for $\Gamma(\frac12)$ substitute $x=\frac{u^2}2$. (b) Induction with Proposition~\ref{prop:gausssum}'s convolution formula.
\end{hint}
\begin{hint}{ex:order}
$U_{(k)}\le t$ iff at least $k$ of the $U_i$ are $\le t$. Differentiate the binomial sum; it telescopes.
\end{hint}
\begin{hint}{ex:cantor}
$C\subseteq C_n$, the union of $2^n$ intervals of length $3^{-n}$. For continuity, show $P(X=x)\le2^{-n}$ for every $n$.
\end{hint}
\begin{hint}{ex:betagamma}
Compute $\mathbb E[\psi(S)\chi(V)]$ for Borel $\psi,\chi\ge0$ as a double integral and change variables in two one-dimensional steps:
$x=\frac{yv}{1-v}$ for fixed $y$, then $y=s(1-v)$ for fixed $v$.
\end{hint}

\section{Solutions}

\begin{solution}{ex:counter}
\idea Mass that escapes to a point is invisible to pointwise limits; the theorems forbid exactly this.
\writeup For $x\le0$, $f_n(x)=0$; for $x>0$, $f_n(x)=0$ once $n>\frac1x$. And $\int f_n=n\cdot\frac1n=1$. The sequence is not monotone ($f_n(x)$ goes
up then drops to $0$), and any $g\ge f_n$ for all $n$ satisfies $g(x)\ge\lfloor\frac1x\rfloor$ on $]0,1[$, which is not integrable since
$\int_{1/(k+1)}^{1/k}\lfloor\frac1x\rfloor\dd x=\frac1{k+1}$ and the harmonic series diverges. Fatou gives $\int\liminf f_n=0<1=\liminf\int f_n$.
\end{solution}

\begin{solution}{ex:densityunique}
\idea Two densities give the same mass to every set, in particular to the set where one beats the other.
\writeup Let $B=\{f>g\}$, a Borel set. Then $\int_Bf=P(X\in B)=\int_Bg$, and both are finite, so $\int(f-g)\mathbf 1_B=0$ with $(f-g)\mathbf 1_B\ge0$. Hence
$(f-g)\mathbf 1_B=0$ a.e., i.e.\ $\lambda(B)=0$. By symmetry $\lambda(\{g>f\})=0$.
\end{solution}

\begin{solution}{ex:tail}
\idea A number is the length of the interval below it; integrate this over $\Omega$ and exchange the order.
\writeup For every $\omega$, $X(\omega)=\int_0^\infty\mathbf 1_{t<X(\omega)}\dd t$. The function $(\omega,t)\mapsto\mathbf 1_{t<X(\omega)}$ is measurable on $\Omega\times[0,\infty[$ (the set
$\{(\omega,t):t<X(\omega)\}$ is the preimage of an open set by the measurable map $(\omega,t)\mapsto X(\omega)-t$). By Tonelli,
$\mathbb E[X]=\int_\Omega\int_0^\infty\mathbf 1_{t<X}\dd t\dd P=\int_0^\infty P(X>t)\dd t$. Likewise $X^2=\int_0^\infty2t\mathbf 1_{t<X}\dd t$. For $X$ with values in $\N$,
$P(X>t)=P(X\ge k)$ for $t\in[k-1,k[$, and the formula is $\sum_{k\ge1}P(X\ge k)$, Theorem~\ref{thm:linearity}(d).
\end{solution}

\begin{solution}{ex:cauchy}
\idea In the Box--Muller representation the angle is uniform and independent of the radius; all three claims are about
the angle.
\writeup By Theorem~\ref{thm:boxmuller}, $(X,Y)$ has the law of $(R\cos\Theta,R\sin\Theta)$ with $R=\sqrt{-2\ln U_1}$ and $\Theta=2\pi U_2$ independent, and every
claim is about the law of $(X,Y)$, so we may use this representation. (a) $Y/X=\tan\Theta$, and $\tan(2\pi U_2)$ has the law of $\tan(\pi U)$ for
$U$ uniform ($\tan$ has period $\pi$, and $2U_2$ mod $1$ is uniform), which is $\tan(\pi(U'-\frac12))$ for the uniform variable $U'=U+\frac12$ mod $1$ (period $\pi$ again), Cauchy by the remark after
Proposition~\ref{prop:inverse}. (b) $\mathbb E|C|=\int_\R\frac{|x|}{\pi(1+x^2)}\dd x=\frac2\pi\lim_{A\to\infty}\frac12\ln(1+A^2)=\infty$. (c) $(X/R,Y/R)=(\cos\Theta,\sin\Theta)$ with $\Theta$ uniform,
independent of $R$, which is a function of $U_1$; and $R^2=-2\ln U_1$ is exponential$(\frac12)$, of mean $2$, by the inverse transform.
\end{solution}

\begin{solution}{ex:gamma}
\idea Integration by parts lowers the exponent; the half-integer case is the Gaussian integral in disguise.
\writeup (a) By parts on $[\varepsilon,A]$ and limits, $\Gamma(a+1)=[-x^ae^{-x}]_0^\infty+a\int_0^\infty x^{a-1}e^{-x}\dd x=a\Gamma(a)$. With $\Gamma(1)=1$, induction gives
$\Gamma(n)=(n-1)!$. With $x=\frac{u^2}2$ (Proposition~\ref{prop:cov}), $\Gamma(\frac12)=\int_0^\infty x^{-1/2}e^{-x}\dd x=\int_0^\infty\frac{\sqrt2}ue^{-u^2/2}u\dd u=\sqrt2\cdot\frac{\sqrt{2\pi}}2=\sqrt\pi$.
(b) If $S$ is Gamma$(n)$ and $E$ exponential$(1)$, independent, then $S+E$ has density, for $z>0$,
$\int_0^z\frac{(z-y)^{n-1}}{(n-1)!}e^{-(z-y)}e^{-y}\dd y=e^{-z}\frac{z^n}{n!}$, the Gamma$(n+1)$ density. Induct from Gamma$(1)$, the exponential law.
\end{solution}

\begin{solution}{ex:order}
\idea The event $U_{(k)}\le t$ is a binomial event; its derivative in $t$ collapses to one term.
\writeup The number $N_t$ of $U_i\le t$ is binomial$(n,t)$, and $U_{(k)}\le t$ iff $N_t\ge k$. So $F(t)=\sum_{j\ge k}\binom njt^j(1-t)^{n-j}$, a polynomial. Using
$j\binom nj=n\binom{n-1}{j-1}$ and $(n-j)\binom nj=n\binom{n-1}j$,
\[
\begin{aligned}F'(t)&=n\sum_{j\ge k}\binom{n-1}{j-1}t^{j-1}(1-t)^{n-j}\eqbr-n\sum_{j\ge k}\binom{n-1}jt^j(1-t)^{n-1-j},\end{aligned}
\]
and the second sum is the first without its $j=k$ term (shift $j\to j+1$). So $F'(t)=n\binom{n-1}{k-1}t^{k-1}(1-t)^{n-k}$, a density of $U_{(k)}$ since $F$ is a polynomial on $[0,1]$, $0$ before and $1$ after. Then $\mathbb E[U_{(k)}]=n\binom{n-1}{k-1}\int_0^1t^k(1-t)^{n-k}\dd t$. By Exercise~\ref{ex:betagamma} (or $n-k$ integrations by parts) the integral is
$\frac{k!\,(n-k)!}{(n+1)!}$, and the product simplifies to $\frac k{n+1}$. The $n$ points cut $[0,1]$ into $n+1$ gaps of the same mean length.
\end{solution}

\begin{solution}{ex:cantor}
\idea The Cantor set is tiny for length, but it is all that the digits $0$, $2$ can reach.
\writeup $C\subseteq C_n$, the set of numbers whose first $n$ ternary digits can be taken in $\{0,2\}$, a union of $2^n$ closed intervals of length
$3^{-n}$; so $\lambda(C)\le(\frac23)^n\to0$. Every value of $X$ has a ternary expansion with digits $2B_k\in\{0,2\}$, so $X\in C$ always. For a
point $x$, $X=x$ forces the first $n$ digits $2B_1,\dots,2B_n$ to be determined by $x$ (two expansions of $X$ with digits in $\{0,2\}$ that
differ first at place $k$ differ by at least $2\cdot3^{-k}-\sum_{j>k}2\cdot3^{-j}=3^{-k}>0$), so $P(X=x)\le2^{-n}$ for all $n$: $F_X$ has no jumps. If
$X$ had a density $f$, then $1=P(X\in C)=\int_Cf\dd\lambda=0$, since $\lambda(C)=0$. So $F_X$, the \emph{Cantor staircase}, is continuous,
non-decreasing, constant on each removed interval, and still climbs from $0$ to $1$.
\end{solution}

\begin{solution}{ex:betagamma}
\idea A pair of independent Gammas, written in ``total and proportion'' coordinates, has a product density; the
Jacobian of $(s,v)\mapsto(sv,s(1-v))$ is $s$, and two one-dimensional changes reproduce it.
\writeup Let $c=\frac1{\Gamma(a)\Gamma(b)}$ and $\psi,\chi\ge0$ Borel on $]0,\infty[$ and $]0,1[$. By Proposition~\ref{prop:joint} and Tonelli,
\[
\begin{aligned}\mathbb E[\psi(S)\chi(V)]&=c\int_0^\infty\!\!\int_0^\infty\psi(x+y)\chi\Bigl(\frac x{x+y}\Bigr)\eqbr\times x^{a-1}y^{b-1}e^{-x-y}\dd x\dd y.\end{aligned}
\]
For fixed $y$, put $x=\frac{yv}{1-v}$, a $C^1$ increasing bijection $]0,1[\to\,]0,\infty[$ with $\frac{\dd x}{\dd v}=\frac y{(1-v)^2}$; then $x+y=\frac y{1-v}$ and $\frac x{x+y}=v$. Exchange the
integrals and, for fixed $v$, put $y=s(1-v)$, so that $x=sv$, $x+y=s$, and the two factors $\frac y{(1-v)^2}$ and $\frac{\dd y}{\dd s}=1-v$ multiply to $s$:
\[
\begin{aligned}\mathbb E[\psi(S)\chi(V)]&=c\int_0^1\!\!\int_0^\infty\psi(s)\chi(v)\eqbr\times(sv)^{a-1}(s-sv)^{b-1}e^{-s}\,s\dd s\dd v\\
&=\int_0^\infty\psi(s)g(s)\dd s\eqbr\times\int_0^1\chi(v)\beta(v)\dd v,\end{aligned}
\]
where
\[
\begin{gathered}g(s)=\frac{s^{a+b-1}e^{-s}}{\Gamma(a+b)},\qq\beta(v)=\frac{\Gamma(a+b)}{\Gamma(a)\Gamma(b)}v^{a-1}(1-v)^{b-1}.\end{gathered}
\]
The function $g$ integrates to $1$ by definition of $\Gamma(a+b)$. With $\psi=\chi=1$ the left side
is $1$, so $\int_0^1\beta=1$: this is the Beta integral, and $\beta$ is a density. Now take $\psi=\mathbf 1_{]0,s]}$ and $\chi=\mathbf 1_{]0,w]}$:
$P(S\le s,V\le w)=\int_0^sg\cdot\int_0^w\beta$. Letting $w\to1$ and then $s\to\infty$ shows that $S$ and $V$ have the densities $g$ and $\beta$, and then
the product form gives their independence (Proposition~\ref{prop:indepcdf}(a)).
\end{solution}

\section*{Sources and further reading}

\begin{itemize}
\item D.~Williams, \emph{Probability with Martingales}, Cambridge University Press, 1991, Chapters~1--6 and~A. Measure
theory for probabilists in a hundred pages, with the $\pi$-$\lambda$ theorem and monotone convergence as the two pillars.
\item H.~Lebesgue, \emph{Sur une g\'en\'eralisation de l'int\'egrale d\'efinie}, Comptes Rendus 132 (1901), 1025--1028. Six pages
that changed analysis.
\item G.~Vitali, \emph{Sul problema della misura dei gruppi di punti di una retta}, Bologna, 1905; S.~Wagon, \emph{The
Banach--Tarski Paradox}, Cambridge University Press, 1985, for the three-dimensional story.
\item A.~N.~Kolmogorov, \emph{Grundbegriffe der Wahrscheinlichkeitsrechnung}, Springer, 1933 (English: \emph{Foundations of
the Theory of Probability}, Chelsea, 1956). Probability as measure theory, in sixty pages.
\item G.~E.~P.~Box and M.~E.~Muller, \emph{A note on the generation of random normal deviates}, Annals of
Mathematical Statistics 29 (1958), 610--611.
\item L.~Devroye, \emph{Non-Uniform Random Variate Generation}, Springer, 1986 (free online). Everything about turning
uniform numbers into other laws.
\item E.~B.~Dynkin, \emph{Foundations of the Theory of Markov Processes}, Pergamon, 1960, where the $\pi$-$\lambda$ argument
became a standard tool.
\end{itemize}
